\lecture{4}
\title{Context-Free Grammars and Pushdown Automata}

\begin{document}

\subsection{Context-Free Grammars: URLs}

Recall from
\href{https://web.mit.edu/6.102/www/sp26/classes/12-grammars-parsing}{Lecture 12}
of 6.1020 (Software Construction)
that the following is a \textbf{context-free grammar (CFG)}.
\begin{verbatim}
url ::= 'http://' hostname (':' port)? '/' 
hostname ::= word '.' hostname | word '.' word
port ::= [0-9]+
word ::= [a-z]+
\end{verbatim}
We might describe the components of this CFG like so:
\begin{itemize}
    \item The \textbf{variables} $V$ are \str{url}, \str{hostname},
    \str{word}, and \str{port}.
    \item The \textbf{terminals} $\Sigma$
    are the digits, the lowercase letters,
    and the symbols $\{\str{.}, \ \str{:}, \ \str{/}\}$.
    \item The \textbf{rules} $R$ are the four lines of the grammar.
    \begin{itemize}
        \item If one more formally says that a rule is
        $(A, w) \in V \times (V \cup \Sigma)^*$, then
        some lines are technically multiple rules.
        \item For example,
        $(\str{hostname}, \str{word}.\str{word}) \in R$
        is one such rule.
    \end{itemize}
    \item The \textbf{start variable} $S \in V$ is \str{url}.
\end{itemize}
The \textbf{language} corresponding to this CFG is the set of valid URLs.

\textbf{Definition.}
A language $A = L(G)$ is a \textbf{context-free language (CFL)}
if it is the language of some CFG $G$.

By replacing all variables with their productions,
we can also describe our URL language
in a single regular expression.
\begin{verbatim}
url ::= 'http://' ([a-z]+ '.')+ [a-z]+ (':' [0-9]+)? '/' 
\end{verbatim}
Note that this cannot be done for all CFGs.
\begin{verbatim}
start ::= '0' start '1' | ''        // cannot be done
start ::= '0' start | ''            // can be done
\end{verbatim}
The problem with the former of these two CFGs is \emph{self-embedding}.
In fact, it turns out that any CFG without self-embedding
can be reduced to a regular expression.
(This fact was not discussed in lecture; proof omitted.)

\subsection{Context-Free Grammars: PEMDAS}

Consider this example of a CFG $G_1 := (V, \Sigma, R, S)$.
\begin{itemize}
    \item Take the variables $V := \{E, T, F\}$.
    \item Take the terminals $\Sigma := \{+, \times, (, ), a\}$.
    \item Take the rules $R := \left\{
        (E, ~ E + T \mid T), ~~
        (T, ~ T \times F \mid F), ~~
        (F, ~ (E) \mid a)
    \right\}$.
    \begin{itemize}
        \item Here, $\mid$ denotes ``or''.
        More formally, $R$ should contain $6$ ordered pairs,
        such as $(E, E+T)$ and $(E, T)$, say.
    \end{itemize}
    \item Take the start variable $S := E$.
\end{itemize}
Then the language $L(G_1)$ of $G_1$
contains words such as $a + a \times a$,
and $(a + a) \times a$, and $a + (a \times a)$.

Each of these words has its own \textbf{parse tree}
that describes its \textbf{derivation} from variables to terminals.

\begin{center}
    \frame{\includegraphics[width=0.25\textwidth]{lecture-04/parse-tree}}
\end{center}

Notice how the structure of the parse tree, itself,
encodes the fact that $\times$ comes before $+$.

\subsection{Context-Free Grammars: Ambiguous PEMDAS}

Now consider the following CFG $G_2 := (V, \Sigma, R, S)$
with the same language.
\begin{itemize}
    \item Take the variables $V := \{E\}$.
    \item Take the terminals $\Sigma := \{+, \times, (, ), a\}$.
    \item Take the rules $R := \left\{
        (E, ~ E + E \mid E \times E \mid (E) \mid a)
    \right\}$.
    \item Take the start variable $S := E$.
\end{itemize}
Then $L(G_2)$ also contains words such as $a + a \times a$,
and $(a + a) \times a$, and $a + (a \times a)$.

But unlike $G_1$, the grammar $G_2$ is \textbf{ambiguous},
in that some words can have multiple \textbf{derivations}
with distinct \textbf{parse trees}.

\begin{center}
    \frame{\includegraphics[width=0.4\textwidth]{lecture-04/ambiguity}}
\end{center}

The ambiguity of this grammar means that
this grammar does not inherently encode the fact that
$\times$ comes before $+$.

\subsection{Pushdown Automata}

Informally, a pushdown automaton is just an NFA
that is granted the additional ability to read and write memory to a stack.

\begin{center}
    \frame{\includegraphics[width=0.5\textwidth]{lecture-04/PDA}}
\end{center}

\textbf{Definition.}
A pushdown automaton (PDA) is a 6-tuple
$P = (Q, \Sigma, \Gamma, \delta, q_0, F)$ where:
\begin{itemize}
    \item The usual NFA components have the same meanings as before.
    \begin{itemize}
        \item The set of \textbf{states} is $Q$.
        \item The \textbf{input alphabet}
        (the set of atoms accepted as input) is $\Sigma$.
        \item The \textbf{start state} is $q_0 \in Q$.
        \item The set of \textbf{accept states} is $F \subseteq Q$.
    \end{itemize}
    \item The \textbf{stack alphabet}
    (the set of atoms accepted by the stack) is $\Gamma$.
    \item The \textbf{transition function} has type
    $\delta \colon Q \times \Sigma_{\epsilon} \times \Gamma_{\epsilon}
    \to \mathcal{P}(Q \times \Gamma_{\epsilon})$
    where $\Sigma_{\epsilon} := \Sigma \cup \{\epsilon\}$
    and $\Gamma_{\epsilon} := \Gamma \cup \{\epsilon\}$.
\end{itemize}

\begin{problem}
    Design a pushdown automaton that recognizes the (nonregular!) language
    $A := \{\str{0}^k\str{1}^k \mid k \geq 0\}$.
\end{problem}

\begin{solution}
    The algorithm, informally, works like so:
    \begin{enumerate}
        \item So long as the input is \str{0},
        push \str{X} onto the stack.
        \item The first time the input is \str{1},
        transition from the start state $q_0$ to the state $q_1$
        without consuming that \str{1}.
        \item So long as the input is \str{1},
        pop \str{X} from the stack.
        \begin{itemize}
            \item If at any point a \str{0} is read as input,
            move to a permanent \str{REJECT} state.
            \item If at any point a \str{1} is read while the stack is empty,
            move to a permanent \str{REJECT} state.
        \end{itemize}
        \item At the end of the input string,
        move to an \str{ACCEPT} state if the stack is empty;
        otherwise, \str{REJECT}.
    \end{enumerate}
    Technically, the hardware of a PDA
    cannot explicitly ``know'' whether the stack is empty or not.
    Here's an easy fix:
    as a Step 0, push \str{END} onto the stack;
    if the top of the stack is ever \str{END},
    the stack must be effectively empty.
    ~ $\blacksquare$
\end{solution}

\begin{problem}
    Design a pushdown automaton that recognizes the (nonregular!) language
    $A := \{ww^R \mid w \in \{\str{0}, \str{1}\}^*\}$.
\end{problem}

\begin{solution}
    If we knew that $w$ had length $10$ (say),
    then all we would do is push the first $10$ input symbols onto the stack,
    then pop the stack $10$ times and compare the result of each pop
    against the next $10$ input symbols read.

    Of course, we don't know the length of $w$ before we read it.
    But the \emph{nondeterminism} of the NFA component of the PDA
    allows us to work around this in the obvious way.
    ~ $\blacksquare$
\end{solution}

\subsection{CFLs are PDA Languages}

We'll finish lecture with a theorem that says
exactly what you'd expect.

\begin{theorem}{CFL $\implies$ PDA}
    If $A$ is a context-free language, then it is recognized by some PDA.
\end{theorem}

\begin{proof}
    Given any $w \in A$, the PDA just uses the stack as scratch work
    to nondeterministically try to guess the derivation of $w$.
    \begin{enumerate}
        \item Initially, the stack contains
        the start variable of the CFG atop \str{END}.
        \item Repeat the following until \str{ACCEPT}
        or \str{REJECT} is reached.  Suppose the top of the stack is $X$.
        \begin{itemize}
            \item If $X$ is a variable,
            then (nondeterministically) replace $X$ with
            one of its productions.
            \item If $X$ is a terminal symbol,
            then pop $X$ from the stack, and
            check if $X$ matches the input.
            \begin{itemize}
                \item If it does match,
                then proceed to the next atom in the input.
                \item If it doesn't match, then \str{REJECT}.
            \end{itemize}
            \item If $X = \str{END}$, then
            \str{ACCEPT} if the input has been entirely read,
            and \str{REJECT} otherwise.
        \end{itemize}
    \end{enumerate}
\end{proof}

\begin{remark}
    Interestingly, the states of the PDA are effectively unused;
    they only exist to move to \str{REJECT} and \str{ACCEPT}.
\end{remark}

It turns out the converse is also true.  In particular,

\begin{theorem}{CFL $=$ PDA}
    A language is context-free if and only if
    it is recognized by some PDA.
\end{theorem}

The proof of the converse is in the textbook,
but we won't go over it in lecture.

\end{document}